\documentclass[11pt]{article}
\usepackage[utf8]{inputenc}
\usepackage{amsmath,amssymb}
\usepackage{hyperref}
\title{Efficient Metal Organic Framework Gas Storage: A Resource-Aware Approach}
\author{Anonymous}
\date{}
\begin{document}
\maketitle

\begin{abstract}
This paper studies Metal Organic Framework Gas Storage. Grounded in a real, citable literature \cite{ref1,ref2,ref3,ref4,ref5,ref6,ref7,ref8},
we propose a focused method and report \textbf{illustrative} experiments.
All references below are real and verifiable (OpenAlex/arXiv); the reported
numbers are simulated to illustrate intended behaviour.
\end{abstract}

\section{Introduction}
Metal Organic Framework Gas Storage is an active research area. This work targets a concrete gap identified from
the recent literature and contributes a method together with an evaluation protocol.

\section{Related Work (real literature)}
The following works, retrieved from public bibliographic APIs, frame our study:
\begin{itemize}
\item Furukawa et al. (2013) — "The Chemistry and Applications of Metal-Organic Frameworks", Science (cited 16741).
\item Li et al. (1999) — "Design and synthesis of an exceptionally stable and highly porous metal-organic framework", Nature (cited 8440).
\item Hafizovic et al. (2008) — "A New Zirconium Inorganic Building Brick Forming Metal Organic Frameworks with Exceptional Stability", Journal of the American Chemical Society (cited 7242).
\item D’Alessandro et al. (2010) — "Carbon Dioxide Capture: Prospects for New Materials", Angewandte Chemie International Edition (cited 4005).
\item Furukawa et al. (2010) — "Ultrahigh Porosity in Metal-Organic Frameworks", Science (cited 3868).
\item Yuan et al. (2018) — "Stable Metal–Organic Frameworks: Design, Synthesis, and Applications", Advanced Materials (cited 2966).
\end{itemize}

\section{Method}
We reduce the dominant computational cost via adaptive resource allocation, activating only the components needed per input. We instantiate this idea for Metal Organic Framework Gas Storage and analyse its cost/quality trade-off.

\section{Experiments (Illustrative)}
\begin{center}
\begin{tabular}{lcc}
System & Primary Metric & Efficiency \\ \hline
Strong baseline & 0.812 & 1.00$\times$ \\
Ablation & 0.828 & 0.92$\times$ \\
\textbf{Ours} & \textbf{0.851} & \textbf{0.71$\times$}
\end{tabular}
\end{center}
\emph{Metrics simulated to illustrate the intended effect; not measured.}

\section{Conclusion}
We connected a real literature to a concrete method for Metal Organic Framework Gas Storage. Future work will
validate the illustrative results with real experiments.

\begin{thebibliography}{9}
\bibitem{ref1} Hiroyasu Furukawa, Kyle E. Cordova, M. O’Keeffe et al.. The Chemistry and Applications of Metal-Organic Frameworks. \emph{Science}, 2013. DOI:10.1126/science.1230444
\bibitem{ref2} Hailian Li, Mohamed Eddaoudi, M. O’Keeffe et al.. Design and synthesis of an exceptionally stable and highly porous metal-organic framework. \emph{Nature}, 1999. DOI:10.1038/46248
\bibitem{ref3} J. Hafizovic, Søren Jakobsen, Unni Olsbye et al.. A New Zirconium Inorganic Building Brick Forming Metal Organic Frameworks with Exceptional Stability. \emph{Journal of the American Chemical Society}, 2008. DOI:10.1021/ja8057953
\bibitem{ref4} Deanna M. D’Alessandro, Berend Smit, Jeffrey R. Long. Carbon Dioxide Capture: Prospects for New Materials. \emph{Angewandte Chemie International Edition}, 2010. DOI:10.1002/anie.201000431
\bibitem{ref5} Hiroyasu Furukawa, Nakeun Ko, Yong Bok Go et al.. Ultrahigh Porosity in Metal-Organic Frameworks. \emph{Science}, 2010. DOI:10.1126/science.1192160
\bibitem{ref6} Shuai Yuan, Liang Feng, Kecheng Wang et al.. Stable Metal–Organic Frameworks: Design, Synthesis, and Applications. \emph{Advanced Materials}, 2018. DOI:10.1002/adma.201704303
\bibitem{ref7} Andrew R. Millward, Omar M. Yaghi. Metal−Organic Frameworks with Exceptionally High Capacity for Storage of Carbon Dioxide at Room Temperature. \emph{Journal of the American Chemical Society}, 2005. DOI:10.1021/ja0570032
\bibitem{ref8} Materials Design and Discovery Group, Hee K. Chae, Diana Y. Siberio-Pérez et al.. A route to high surface area, porosity and inclusion of large molecules in crystals. \emph{Nature}, 2004. DOI:10.1038/nature02311
\end{thebibliography}
\end{document}
